\chapter{End-to-end ROS pipelines}
\label{ch:pipelines}

This chapter walks the live system as pipelines. Each arrow is a ROS topic
(or Gazebo Transport channel). Chapter~\ref{ch:ros-catalog} then expands every
node.

\section{Pipeline A --- Motion control (closed loop)}
\begin{verbatim}
Gazebo dynamics
   |  (Gazebo Transport: /world/.../dynamic_pose/info)
   v
gazebo_pose2d  --->  /asvK/pose2d
   |
   |  also: ros_gz_bridge ---> /asvK/odom
   v
los_path_follower
   subs: /asvK/odom , /asvK/plan   (plan is TRANSIENT_LOCAL)
   pubs: /asvK/cmd_vel , /asvK/plan/active , /asvK/nav/debug
   v
thrust_mixer
   subs: /asvK/cmd_vel
   pubs: /model/<boat>/joint/left|right_propeller_joint/cmd_thrust
   v
ros_gz_bridge (ROS_TO_GZ) ---> Gazebo propeller joints
\end{verbatim}

\paragraph{Who owns the path?}
\texttt{mission\_manager} (or \texttt{spiral\_demo}) publishes
\texttt{/asv\(K\)/plan}. The LOS follower latches the latest Path (QoS
transient local) and tracks it until a new plan arrives. During HOLD / COMPLETE
the mission manager publishes a trivial hold path near the current pose.

\paragraph{Logic order inside one control cycle.}
\begin{enumerate}[leftmargin=*]
  \item Odometry updates boat pose for LOS geometry.
  \item Cross-track error and lookahead yield a commanded heading.
  \item Yaw PID $\rightarrow$ \texttt{angular.z}; surge policy $\rightarrow$
        \texttt{linear.x}.
  \item Mixer maps \((v,\omega)\) to left/right thrust floats.
  \item Bridge writes thrusts into Gazebo; hydrodynamics integrate.
\end{enumerate}

\section{Pipeline B --- Magnetometry to cleaned anomaly}
\begin{verbatim}
Gazebo magnetometer
   v
ros_gz_bridge ---> /asvK/mag/gazebo   (sensor_msgs/MagneticField, Tesla)
   v
mag_driver
   also uses: /asvK/pose2d , /asvK/cmd_vel
   optionally overlays planted dipole A/(r^3+s^3)
   pubs: /asvK/mag/raw                 (boat_msgs/MagReading, nT)
   v
mag_filter
   spike reject + moving-average LPF
   pubs: /asvK/mag/filtered , /asvK/mag/filter_status
   v
calibration_node
   spatial heading-binned baseline + temporal high-pass
   pubs: /asvK/mag/anomaly             (boat_msgs/MagAnomaly)
         /asvK/calibration/status      (CALIBRATING|READY|DEGRADED)
\end{verbatim}

\paragraph{Why three stages?}
\begin{itemize}[leftmargin=*]
  \item \texttt{mag\_driver} owns units, pose stamping, motor flag, and the
        synthetic dipole plant used only in simulation evaluation.
  \item \texttt{mag\_filter} owns robust denoising so calibration does not
        learn spikes.
  \item \texttt{calibration\_node} owns ambient baseline learning and the
        cleaned anomaly feature consumed by Bayes and verification.
\end{itemize}

\section{Pipeline C --- Shared belief and localization}
\begin{verbatim}
/asv1/mag/anomaly  --+
                     +-->  bayes_fusion  -->  /swarm/belief/map
/asv2/mag/anomaly  --+                    -->  /swarm/belief/peak
                                          -->  /swarm/belief/peak_probability
                                          -->  /swarm/belief/centroid
                                          -->  /swarm/belief/centroid_mass
                                          -->  /swarm/belief/centroid_spread
                                          -->  /swarm/belief/fix
                                          -->  /swarm/belief/fix_rms
                                          -->  /swarm/belief/fix_samples
\end{verbatim}

\paragraph{Internal order in \texttt{bayes\_fusion}.}
\begin{enumerate}[leftmargin=*]
  \item Receive \texttt{MagAnomaly}; ignore if not calibrated (ABSTAIN).
  \item Classify HIT / MISS / ABSTAIN against thresholds.
  \item Update every grid cell with \(P_{\mathrm{det}}(d)\) likelihood; renormalize.
  \item Compute discrete peak and belief-weighted centroid (+ spread).
  \item Buffer strong samples into dipole fitter; optionally publish continuous fix.
\end{enumerate}

\section{Pipeline D --- Mission behaviour and swarm verify}
\begin{verbatim}
mission_manager (/asvK)
   subs: pose2d, mag/anomaly, calibration/status,
         /swarm/belief/{map,peak,peak_probability,centroid},
         /swarm/verify/request , /swarm/mission/halt
   pubs: /asvK/plan , /asvK/mission/state , /asvK/mission/mode ,
         /asvK/mission/info_gain ,
         /swarm/verify/declare , /swarm/verify/result

verify_coordinator (global, once)
   subs: /swarm/verify/declare , /swarm/verify/result
   pubs: /swarm/verify/request , /swarm/mission/complete ,
         /swarm/mission/halt , /swarm/mission/status
\end{verbatim}

\paragraph{Dual-ASV happy path.}
\begin{enumerate}[leftmargin=*]
  \item Both boats lawnmower Voronoi regions (\texttt{GLOBAL\_SEARCH}).
  \item Belief HITs raise \(p^\star\); nearby boat enters
        \texttt{TARGET\_SEARCH} (MI then spiral).
  \item Discoverer declares candidate on \texttt{/swarm/verify/declare} and
        HOLDs.
  \item Coordinator assigns the other ASV on \texttt{/swarm/verify/request}.
  \item Verifier orbits, counts strong confirmations, publishes
        \texttt{VerifyResult}.
  \item Coordinator latches COMPLETE / HALT; plotter banner shows
        \textbf{CONFIRMED}.
\end{enumerate}

\section{Pipeline E --- Operator visualization}
\texttt{trajectory\_plotter} (global in multi-ASV mode) subscribes to every
boat's \texttt{pose2d}, \texttt{plan/active}, \texttt{mag/anomaly},
\texttt{mission/state}, plus all \texttt{/swarm/belief/*} and verify/complete
topics. It does not publish commands. The bottom status strip mirrors mission
phase chips: SEARCH $\rightarrow$ HUNT $\rightarrow$ VERIFY $\rightarrow$
CONFIRMED.

\section{Minimum viable process set}
To move one boat under LOS alone you need: Gazebo, bridges,
\texttt{gazebo\_pose2d}, \texttt{los\_path\_follower}, \texttt{thrust\_mixer},
and someone publishing a Path. To run the full magnetic mission you additionally
need the sensing trio, \texttt{bayes\_fusion}, \texttt{mission\_manager}, and
(for cooperative verify) \texttt{verify\_coordinator}.
